% Generated draft: requires mathematical and source review.
\documentclass[11pt]{article}
\usepackage[margin=25mm]{geometry}
\usepackage{fontspec}
\setmainfont{Latin Modern Roman}
\usepackage{amsmath,amssymb,mathtools,unicode-math}
\setmathfont{Latin Modern Math}
\usepackage{enumitem,needspace,xurl,hyperref}
\setlength{\parindent}{0pt}
\setlength{\parskip}{5pt}
\setlength{\emergencystretch}{3em}
\newcommand{\N}{\mathbb{N}}
\newcommand{\Z}{\mathbb{Z}}
\newcommand{\Q}{\mathbb{Q}}
\newcommand{\R}{\mathbb{R}}
\newcommand{\C}{\mathbb{C}}
\newcommand{\F}{\mathbb{F}}
\newcommand{\A}{\mathbb{A}}
\newcommand{\PP}{\mathbb{P}}
\newcommand{\E}{\mathbb{E}}
\newcommand{\eps}{\varepsilon}
\newcommand{\dd}{\,\mathrm{d}}
\begin{document}
\section*{UnsolvedMath ID: 30005909}
\textbf{Source problem code:} catalogue entry
\subsection{Definitions and mathematical statement}
For an object \(S\) for which the relevant products and quotient stacks are defined, the symmetric group \(\mathfrak S_n\) acts on \(S^n\) by permuting the factors. The corresponding symmetric quotient stack is
\[
\mathcal X_n=[S^n/\mathfrak S_n].
\]
This is a conventional interpretation of ``symmetric quotient stacks''; the supplied record does not specify the permitted objects \(S\), the base field, or the range of \(n\).

An ordinary Heisenberg Lie algebra associated with a vector space \(V\) over a field \(K\) and an alternating bilinear form \(\omega\) is the central extension
\[
\mathfrak h(V,\omega)=V\oplus Kc,\qquad
[v,w]=\omega(v,w)c,\qquad [c,v]=0.
\]
An action on a vector space \(M\) means a Lie algebra homomorphism \(\mathfrak h(V,\omega)\to\operatorname{End}_K(M)\), where the target has the commutator bracket. These definitions explain ordinary Heisenberg terminology only: the source does not identify \(K\), \(V\), \(\omega\), the central action, or whether a graded or super version is intended.

\textbf{Source question.} Is there a Heisenberg algebra controlling the Hochschild--Serre cohomology of symmetric quotient stacks?

\textbf{Unresolved formulation.} The supplied record does not define ``Hochschild--Serre cohomology'' in this setting or the requirement expressed by ``controlling.'' In particular, it does not specify whether the intended structure is an action on the cohomology of each \(\mathcal X_n\), an action on a sum over different \(n\), or a stronger description by a Fock-space representation. Consequently, the question above is the preserved research target, but the evidence does not determine a fully quantified existence statement or prescribed commutation relations.

\textbf{Retained contextual qualifications.} The appended source excerpt concerns a space denoted \(S\), its Hilbert schemes written \(\operatorname{Hilb}^n S\), negative Hochschild--Serre cohomology of \(S\), and a calculation of \(\mathrm{HH}^{2}(\operatorname{Hilb}^n S)\). It says that, to prove an unavailable assertion ``(2),'' one takes \(k=0\), uses all negative Hochschild--Serre cohomology of \(S\), and subsequently takes \(\ast=2\). It further describes extracting geometric deformations in a Hochschild--Kostant--Rosenberg decomposition by bootstrapping from earlier results and cancelling contributions in the Hochschild--Serre calculation. However, the component notation is damaged: it contains \texttt{V2}, \texttt{H2 (Hilbn S, OHilbn S )}, and \texttt{H2 (Hilbn S, THilbn S )}. The definitions of \(k\), the grading denoted by \(\ast\), assertion (2), the earlier results, the component notation, and the hypotheses on \(S\) are absent. These fragments therefore supply context without establishing additional hypotheses or a precise auxiliary assertion.
\subsection{Short English statement}
The question asks whether a suitable Heisenberg algebra governs Hochschild--Serre cohomology for symmetric quotient stacks. The supplied evidence leaves the cohomology theory, geometric hypotheses, representation space, and meaning of ``controlling'' unspecified. This is a statement-only draft requiring human mathematical review.
\subsection{Sources}
\begin{itemize}
\item Supplied local UnsolvedMath registry record, permanent ID \(30005909\), title \emph{Heisenberg Actions on Hochschild--Serre Cohomology}, source problem code \texttt{catalogue entry}, in the \emph{Oberwolfach Reports Open Problems} collection. The record supplies the year 2024 and no named proposer. Its status is labelled \texttt{open}; its statement status is labelled \texttt{corrected\_verified}, with the supplied qualification that only formula delimiters, notation, grammar, and typography were repaired and that mathematical content and scope were preserved. These are registry labels, not independent verification here.
\item Supplied report citation: \emph{Hochschild (Co)Homology and Applications} (2024), \emph{Oberwolfach Reports}, pp. 1097--1162, DOI 10.4171/owr/2024/20. \url{https://doi.org/10.4171/owr/2024/20}. The record describes Oberwolfach Reports as documenting mathematical workshops and their open-problem discussions. No precise page, speaker attribution, or complete surrounding mathematical setup is supplied for this question.
\item Belmans, Fu and Krug, \emph{Hochschild cohomology of Hilbert schemes}. Supplied literature reference: \url{https://arxiv.org/abs/2309.06244}.
\item Anno, Baranovsky and Logvinenko, \emph{Hochschild homology}. Supplied literature reference: \url{https://arxiv.org/abs/2512.25039}.
\item Supplied literature assessment, checked on exactly \texttt{2026-08-22}: the record labels the problem open and states that Heisenberg/Fock-space structures are known for related Grothendieck and Hochschild-homology constructions, while the requested structure for Hochschild--Serre cohomology is not established. This is supplied dated curation derived from the OWR curation, not an independently verified literature assessment. The record explicitly calls for independent verification of time-sensitive status and cited claims; no source visit or such verification is claimed here.
\end{itemize}
\end{document}
